\chapter{The Complete Algorithm}
\label{ch:fullalgo}

\begin{objectives}
\begin{itemize}
  \item Assemble every component into one event-driven estimator, with the correct order of operations.
  \item Estimate the computational budget and map the algorithm onto the project's software architecture.
  \item Diagnose the most common failures from their symptoms.
\end{itemize}
\end{objectives}

% ------------------------------------------------------------------------------
\section{The event loop}

The estimator consumes a single, time-ordered stream of sensor events. Milestone 3 builds it by merging the presorted IMU,
image (and, for testing, pose) arrays with a two-pointer merge, with the rule that an IMU sample precedes any other event
with the same timestamp. \Cref{alg:msckf} shows the loop.

\begin{algorithm}[t]
\caption{Monocular MSCKF}
\label{alg:msckf}
\begin{algorithmic}[1]
\Require event stream, calibration, noise parameters
\LineComment{load phase: allocation allowed}
\State parse data and calibration; build the event arrays; construct the filter and front end with fixed-max workspaces
\State initialize $(\hat x, P)$ (ground truth in the project, logged; static initialization otherwise)
\LineComment{streaming phase: no allocation in the core}
\For{each event $e$ in timestamp order}
  \If{$e$ is an IMU sample at $t$}
    \State propagate nominal state and covariance from $t_{prev}$ to $t$ \Comment{\cref{ch:strapdown,ch:errordyn}}
  \ElsIf{$e$ is an image at $t_c$}
    \State propagate to $t_c$ (interpolated IMU) \Comment{exact image time}
    \State front end: track, verify, undistort; update track histories \Comment{\cref{ch:frontend}}
    \State process tracks that were lost in this image \Comment{\cref{alg:msckf-update}}
    \If{$N = $ \code{kMaxClones}}
      \State process tracks with an observation at the oldest clone
      \State marginalize the oldest clone in place \Comment{\cref{ch:cloning}}
    \EndIf
    \State clone $(\hat p, \hat q)$; $P \gets JPJ\T$; attach this image's observations to the new clone
  \EndIf
  \State log state, $\diag(P)$, NIS and gate statistics to the preallocated buffer
\EndFor
\end{algorithmic}
\end{algorithm}

Three ordering rules are essential.
\begin{enumerate}
  \item \textbf{Propagate before you update.} A measurement at time $t$ must be applied to the state at time $t$, never to the
        state at the previous IMU sample.
  \item \textbf{Use a clone's observations before you delete it.} Every feature with an observation at the oldest clone is
        processed before that clone is marginalized; otherwise its information is lost.
  \item \textbf{Use each observation once.} After a track has been used, its observations are discarded. If the feature keeps
        being tracked, later observations form a new track.
\end{enumerate}
With the order of \cref{alg:msckf}, the window never holds more than \code{kMaxClones} clones, which keeps the covariance within its
fixed maximum size. A feature processed because its oldest clone is leaving cannot use the current image's observation in that
update, because the new clone does not exist yet. Implementations that clone first need room for one extra clone. Either
choice is valid; the storage bound must match it.
The planar filter of \cref{ch:planar} runs exactly this loop (\cref{alg:planar-msckf}) in a few lines of Python, which makes it a
convenient place to experiment with the ordering rules before committing to them in C++.

\begin{figure}[t]
\centering
\begin{tikzpicture}[xscale=0.9, yscale=1.0]
  \draw[->, thick] (0,0) -- (15.5,0) node[right, font=\footnotesize] {$t$};
  \foreach \x in {0,0.25,...,15.0} \draw[black!45] (\x,-0.08) -- (\x,0.08);
  \node[font=\scriptsize, black!60, below] at (0.6,-0.15) {IMU, 200 Hz};
  \foreach \k/\x in {1/1.25, 2/3.75, 3/6.25, 4/8.75, 5/11.25, 6/13.75} {
    \draw[primary, thick] (\x,-0.25) -- (\x,0.6);
    \draw[fill=primarylight, draw=primary] (\x-0.25,0.6) rectangle (\x+0.25,1.0);
    \node[font=\scriptsize] at (\x,0.8) {$C_\k$};
  }
  \node[font=\scriptsize, primary, below] at (3.75,-0.25) {image (20 Hz)};
  % features
  \draw[plotgreen, very thick] (1.25,1.35) -- (8.75,1.35) node[right, font=\scriptsize] {track $f_1$ lost $\Rightarrow$ update};
  \fill[plotgreen] (1.25,1.35) circle (2pt) (3.75,1.35) circle (2pt) (6.25,1.35) circle (2pt) (8.75,1.35) circle (2pt);
  \draw[plotorange, very thick] (3.75,1.85) -- (13.75,1.85);
  \fill[plotorange] (3.75,1.85) circle (2pt) (6.25,1.85) circle (2pt) (8.75,1.85) circle (2pt) (11.25,1.85) circle (2pt) (13.75,1.85) circle (2pt);
  \node[font=\scriptsize, plotorange, align=left, right] at (13.85,1.95) {$f_2$: oldest clone\\leaving $\Rightarrow$ update};
  \draw[decorate, decoration={brace, amplitude=4pt}] (3.5,2.25) -- (14.0,2.25) node[midway, above=5pt, font=\scriptsize] {window (example \code{kMaxClones} = 5)};
  \draw[conv, ->] (1.25,2.5) -- (1.25,1.1) node[pos=0, above, font=\scriptsize, align=center] {$C_1$ marginalized\\when $C_6$ arrives};
\end{tikzpicture}
\caption{Timeline of the MSCKF. The IMU propagates at 200~Hz; each image creates a clone; feature tracks are processed when
they end or when the clone holding their oldest observation is about to leave the window.}
\label{fig:timeline}
\end{figure}

% ------------------------------------------------------------------------------
\section{Initialization and the first seconds}

The project initializes from ground truth, which removes the hardest failure mode from early testing and is logged as a use of
ground truth (Milestone 10b). With static initialization (\cref{ch:strapdown}), the first seconds are delicate: the yaw is
arbitrary, the accelerometer bias is entangled with tilt, and the window contains few clones with little parallax. Generous
initial covariances, and a check that enough features with sufficient parallax exist before trusting visual updates, avoid
early divergence. Inspect the IMU data of each EuRoC sequence before running on it: a static start followed by vigorous
motion is the easiest case for initialization, and a fast start the hardest.

% ------------------------------------------------------------------------------
\section{Computational budget}

With $N = 10$ clones the error state has $n = 75$ dimensions. Rough flop counts per operation:
\begin{center}\small
\begin{tabular}{@{}l l r@{}}
\toprule
Operation & Frequency & Cost (flops) \\
\midrule
Nominal propagation & 200 Hz & $\sim10^2$ \\
$P_{II}$, $P_{IC}$ propagation & 200 Hz & $\sim4\times10^4$ \\
Cloning ($P \gets JPJ\T$ as block copies) & 20 Hz & $\sim10^4$ \\
Per feature: triangulation, projection, gating ($M = 6$) & per track & $\sim10^5$ \\
QR compression of $d = 450$ rows & per update & $\sim5\times10^6$ \\
Update ($S$, gain, Joseph) with $n = 75$ & per update & $\sim6\times10^6$ \\
Front end (KLT on 150 features) & 20 Hz & dominant in time \\
\bottomrule
\end{tabular}
\end{center}
The filter core needs on the order of $10^7$ flops per image, a few milliseconds on an embedded processor and well under a
millisecond on a desktop. The front end is usually the bottleneck in wall-clock time. The cubic growth with $n$ is why the window
is short and why QR compression matters for large stacks.

% ------------------------------------------------------------------------------
\section{Software architecture}

\begin{figure}[t]
\centering
\begin{tikzpicture}[node distance=8mm and 10mm]
  \node[block, minimum width=48mm, align=left] (core) {\textbf{vio\_core} (static, Eigen only)\\ \footnotesize so3, imu\_integrator, es\_ekf,\\ \footnotesize camera\_model, triangulation, msckf\_update};
  \node[block, minimum width=34mm, above left=10mm and -14mm of core, align=left] (io) {\textbf{vio\_io}\\ \footnotesize dataset\_loader, telemetry\\ \footnotesize + YAML parser};
  \node[block, minimum width=34mm, above right=10mm and -14mm of core, align=left] (fe) {\textbf{vio\_frontend}\\ \footnotesize visual\_frontend\\ \footnotesize + OpenCV};
  \node[block, minimum width=30mm, above=24mm of core, fill=white] (app) {\textbf{vio\_app}\\ \footnotesize main.cpp};
  \draw[->, thick] (io) -- (core); \draw[->, thick] (fe) -- (core);
  \draw[->, thick] (app) -- (io); \draw[->, thick] (app) -- (fe); \draw[->, thick] (app) -- (core);
  \node[font=\footnotesize, align=center, below=2mm of core] {arrows: ``depends on''. OpenCV never reaches \code{vio\_core};\\ each test links only the library it tests.};
\end{tikzpicture}
\caption{Library targets of the project (AGENTS.md). Dependencies point one way, so the estimator core can be built, tested and
reasoned about without the I/O and vision libraries.}
\label{fig:architecture}
\end{figure}

\Cref{fig:architecture} shows how the algorithm maps to the project's targets. The zero-allocation discipline applies to
\code{vio\_core}, which contains everything executed per IMU sample and per update. The front end may allocate but must reuse its
buffers, and the I/O layer allocates freely during loading.

% ------------------------------------------------------------------------------
\section{Debugging guide}

\begin{table}[t]
\centering\small
\begin{tabularx}{\textwidth}{@{}>{\raggedright\arraybackslash}p{4.4cm} X@{}}
\toprule
\textbf{Symptom} & \textbf{Likely causes} \\
\midrule
Altitude diverges at about $2g$ in IMU-only propagation & gravity or specific-force sign; $g_w$ added twice \\
Position diverges within seconds even with ground-truth init & $R_{wb}$ used as $R_{bw}$; quaternion $w/z$ swapped at load; gyro and accel columns or biases swapped; degrees vs radians \\
Rotation matrices slowly scale vectors & missing quaternion normalization \\
NaN after a while & missing Taylor branch; $\Log$ near $\pi$; $\Jrinv$ evaluated beyond $\pi$; feature at zero depth; Cholesky of a non-PD $S$ \\
Corrections systematically make things worse & sign of $H$ or of the residual; injection on the wrong side of the quaternion; global vs local convention mixed \\
Covariance loses symmetry or positive definiteness & simple update form instead of Joseph; missing symmetrization; aliasing in Eigen block products \\
Nearly all features fail the gate & overconfident $P$; $T_{bc}$ inverted; pixel noise used in normalized units (or vice versa); distortion not removed; timestamps misaligned \\
Gate rejects nothing & $R_z$ or $\sigma_{uv}$ far too large; gate threshold indexed with the wrong dimension \\
Slowly growing yaw error with healthy NIS & expected EKF inconsistency along unobservable yaw (\cref{ch:observability}) \\
Scale drift & little acceleration (degenerate motion); accelerometer noise or bias badly tuned \\
Works on synthetic data, fails on EuRoC & extrinsics, time offset, distortion, noise inflation, initialization, IMU gaps \\
Results differ between runs & uninitialized memory (run the sanitizers); unseeded random generators \\
\bottomrule
\end{tabularx}
\caption{Symptoms and their most frequent causes.}
\label{tab:debugging}
\end{table}

\Cref{tab:debugging} collects the failures that experience shows are most common. The milestone structure of the project is itself
the most effective debugging tool: each layer is verified against known truth before the next is built on top of it, so a new
failure is almost always in the newest layer.

\section{What to expect on EuRoC}

Published monocular filter-based estimators reach absolute trajectory errors of the order of 0.1 to 0.5~m on the machine-hall
sequences, depending on features, tuning and initialization, with the difficult sequences (fast motion, motion blur, dark
regions) considerably worse~\cite{delmerico2018benchmark}. Milestone 10b sets a first-pass target of 0.5~m on MH\_01 and a stretch target of
0.2~m. Yaw and position drift without bound; roll and pitch remain accurate throughout; velocity and biases converge within the
first seconds of excitation.

\begin{takeaways}
\begin{itemize}
  \item One time-ordered stream; propagate to each event's time before applying it.
  \item At each image: process lost tracks; if the window is full, process the oldest clone's tracks and marginalize it; then clone.
  \item The core costs about $10^7$ flops per image; the front end usually dominates the run time.
  \item Dependencies point one way: \code{vio\_core} knows nothing about files, YAML or OpenCV.
  \item Most failures map to a small set of sign, frame, unit and timing errors; test each layer against truth.
\end{itemize}
\end{takeaways}

\begin{exercises}
  \item Write the event-loop pseudocode for the alternative ordering (clone first, then process and marginalize). What capacity
        does it require, and what does it gain?
  \item Instrument a run to record the number of features processed per update, their observation counts, and the stacked row
        count $d$. Choose the compile-time row cap from the distribution.
  \item For each row of \cref{tab:debugging}, design a unit or integration test that would have caught the cause before the full
        system was assembled.
\end{exercises}
% Solutions to the exercises of Chapter 28.
\begin{solutions}
  \item With the clone created first, an image at $t_c$ is handled as follows:
        \begin{enumerate}[label=\arabic*:, leftmargin=2.2em, itemsep=0pt]
          \item propagate to $t_c$; run the front end;
          \item clone $(\hat p, \hat q)$; $P \gets JPJ\T$; attach this image's observations to the new clone (now $N + 1$ clones);
          \item process the tracks lost in this image;
          \item \textbf{if} $N + 1 > $ \code{kMaxClones}: process the tracks with an observation at the oldest clone (including the current observation), then
                marginalize the oldest clone.
        \end{enumerate}
        The window transiently holds \code{kMaxClones}$\,+\,1$ clones, so every fixed-max buffer ($P$, the stacked $H$, the clone container) must be sized for
        $15 + 6(\text{\code{kMaxClones}} + 1)$ states. The gain: a track processed because its oldest clone is leaving can include the current image's observation. That is
        its newest view and usually the one with the most parallax relative to the oldest, so it adds two rows per track and lengthens the baseline. The price
        is an update on a state six dimensions larger, about $(81/75)^3 \approx 1.26$ times the cost for \code{kMaxClones}$\,= 10$.

  \item Log, for every update, the number of features processed, each feature's $M_j$, and $d = \sum_j(2M_j - 3)$, and plot the distribution over all sequences. The
        planar filter of \cref{ch:planar} ($N_{\max} = 10$, 4\,500 updates) shows the typical shape. On average 2.4 features are processed per update. Most tracks reach the full
        window ($M = 10$ for 90\% of them), because they are processed when their oldest clone leaves. The rows per update have median 16 and 99th percentile 56,
        but the maximum is 175. The tail comes from bursts: when the window first fills, or when many features leave the view together, nearly every visible
        feature is processed at once. The worst case is therefore (visible features)$\,\times\,$(rows per full track): $150\times(2\cdot10 - 3) = 2550$ rows for the
        project. There are two sound ways to set the cap:
        \begin{itemize}
          \item size the workspace for the worst case: $2550\times76$ doubles are 1.5~MB, allocated once at load time;
          \item better, compress incrementally: append features until the next one would not fit, QR-compress the stack to $n$ rows
                (\cref{prop:compression}), and continue. The workspace then needs only $n$ plus one feature's rows ($75 + 17$), and no feature is ever dropped.
        \end{itemize}
        A cap at a percentile of the distribution without either measure would silently discard information in exactly the bursts that matter. If it is used,
        the dropped features must be counted and logged.

  \item One test per row of \cref{tab:debugging}:
        \begin{itemize}
          \item \emph{$2g$ altitude divergence:} propagate a level, stationary IMU ($a_m = (0, 0, 9.81)$, $\omega_m = 0$) for 10~s from rest; position and velocity must stay
                zero to rounding. Then propagate free fall ($a_m = 0$) and require $v_z = -9.81\,t$.
          \item \emph{Divergence within seconds:} loader tests on a fixture row with distinct values in every column (catching $w$/$z$ swaps, gyro/accel and bias column
                order, degrees). A coning test against an analytic attitude. On EuRoC, propagate IMU-only from ground truth for 1~s and compare with ground truth.
          \item \emph{Vectors slowly scaled:} after $10^6$ random multiplicative updates, $\big|\|q\| - 1\big| < 10^{-12}$.
          \item \emph{NaN:} $\Exp$, $\Log$ and $\Jrinv$ at $0$, $10^{-12}$, $\pi - 10^{-9}$ and exactly $\pi$; triangulation must reject zero and negative depths; feeding a non-PD $S$ must return an
                error, not a NaN; a randomized test asserting finite outputs.
          \item \emph{Corrections make things worse:} finite-difference tests of every Jacobian on $\bplus/\bminus$. An update with a noise-free measurement must reduce the
                error. A mutation check: flip a sign in $H$ and confirm the test fails.
          \item \emph{Covariance loses symmetry or definiteness:} after $10^4$ random updates, $\|P - P\T\| < 10^{-12}\|P\|$ and $\lambda_{\min}(P) \ge 0$. Compare in-place block products with an
                out-of-place reference.
          \item \emph{All features gated out:} a synthetic VIO run with the true noise must reject about 5\% of the features (say 3--8\%). A camera-model test projects a hand-computed
                point through $T_{bc}$ from \code{sensor.yaml}. A unit test checks $\sigma_{uv} = \sigma_{px}/f$.
          \item \emph{Gate rejects nothing:} the same run, with a lower bound on the rejection rate; the $\chi^2$ table checked against reference values, indexed by
                $\dim r$.
          \item \emph{Growing yaw error:} per-component NEES on synthetic data; position, roll and pitch must be consistent, and yaw growth is documented as expected.
          \item \emph{Scale drift:} two synthetic runs, with low and high acceleration; the scale error must be bounded in the excited run.
          \item \emph{Synthetic fine, EuRoC fails:} parse the real \code{T\_BS} (row-major) and intrinsics, and compare with hand-copied values; distortion round trip at the corners and
                edges (Exercise~20.5); merge-order test (IMU before image at equal times); IMU-gap detection ($\Dt > 2\times$ nominal flagged).
          \item \emph{Run-to-run differences:} run the pipeline twice in one test and compare outputs bitwise; sanitizers in CI; seeded generators.
        \end{itemize}
\end{solutions}

